\begin{description}
\item[(T11.1)] From the graph, just before the peak of emission, the time
  period of gravitational waves is approximately $(0.007 \pm 0.004)$\,s.
  \[ T_0 \approx 0.007\,\mathrm{s} \]
  \hfill(2 points)

  That is, the frequency of the gravitational waves is
  $f_0 \approx 142.86$\,Hz. \hfill(1 point)

\item[(T11.2)] Since $m \ll M$ then, by Kepler's third law
  \[ f_{\mathrm{orbital}} = \frac{1}{2\pi}\sqrt{\frac{GM}{r^3}} \]
  \hfill(4 points)

  Hence the frequency of the gravitational waves is
  \[ f_{\mathrm{grav}} = 2 f_{\mathrm{orbital}}
     = \frac{1}{\pi}\sqrt{\frac{GM}{r^3}} \]
  \hfill(1 point)

  The frequency will be maximum when $r = R_{\mathrm{MS}}$. \hfill(1 point)

  For main sequence stars,
  \[ \frac{R_{\mathrm{MS}}}{R_\odot}
     = \left(\frac{M_{\mathrm{MS}}}{M_\odot}\right)^{\alpha}
     \qquad\therefore\quad R_{\mathrm{MS}}
     = R_\odot \left(\frac{M_{\mathrm{MS}}}{M_\odot}\right)^{\alpha} \]
  \hfill(1 point)
  \begin{align*}
    \therefore f_{\mathrm{MS}}
      &= \frac{1}{\pi}\sqrt{\frac{G M_{\mathrm{MS}}}{R_\odot^3}}
         \left(\frac{M_\odot}{M_{\mathrm{MS}}}\right)^{3\alpha/2} \\
      &= \frac{1}{\pi}\sqrt{\frac{G M_\odot}{R_\odot^3}}
         \left(\frac{M_\odot}{M_{\mathrm{MS}}}\right)^{(3\alpha-1)/2} \\
    f_{\mathrm{MS}}
      &= \frac{1}{\pi}\sqrt{\frac{G M_\odot}{R_\odot^3}}
         \left(\frac{M_{\mathrm{MS}}}{M_\odot}\right)^{(1-3\alpha)/2}
  \end{align*}
  \hfill(3 points)

\item[(T11.3)] For possible values of $\alpha$ given in the question, the
  exponent $\frac{1-3\alpha}{2}$ is negative. Thus, if
  $M_{\mathrm{MS}} > M_\odot$, the frequency will be smaller. Thus, for highest
  possible frequency coming from a main sequence star, you should take lowest
  possible mass i.e.\ $\alpha = 1.0$. \hfill(4 points)
  \begin{align*}
    f_{\mathrm{MS,max}}
      &= \frac{1}{\pi}\sqrt{\frac{G M_\odot}{R_\odot^3}}
         \left(\frac{M_{\mathrm{MS}}}{M_\odot}\right)^{\left(\frac{1-3\times 1}{2}\right)} \\
      &= \frac{1}{\pi}\sqrt{\frac{G M_\odot}{R_\odot^3}}
         \times \frac{M_\odot}{M_{\mathrm{MS}}}
  \end{align*}
  \hfill(2 points)

  The frequency of gravitational waves will be given by,
  \begin{align*}
    f_{\mathrm{MS,max}}
      &= \frac{1}{\pi}\sqrt{\frac{6.674 \times 10^{-11} \times 1.989 \times 10^{30}}
         {(6.955 \times 10^{8})^{3}}} \times \frac{1}{0.08} \\
    f_{\mathrm{MS,max}} &= 2.5\,\mathrm{mHz}
  \end{align*}
  \hfill(3 points)

\item[(T11.4)] The maximum frequency would be when $r = R_{\mathrm{WD}}$.
  \hfill(1 point)

  We use the notation $R_{WD\odot}$ for the radius of a solar mass white
  dwarf. Then for white dwarfs
  % sic: the official solution prints the mass ratio below as M_WD/M_sun;
  % from R \propto M^{-1/3} it should be M_sun/M_WD, which is what the
  % subsequent lines (and the numerical answer) actually use.
  \[ R_{WD}^3 = R_{WD\odot}^3 \frac{M_{WD}}{M_\odot} \]
  \hfill(1 point)
  \begin{align*}
    f_{\mathrm{WD}}
      &= \frac{1}{\pi}\sqrt{\frac{G M_{\mathrm{WD}}}{R_{WD}^3}} \\
      &= \frac{1}{\pi}\sqrt{\frac{G M_\odot}{R_{WD\odot}^3}}\,
         \frac{M_{WD}}{M_\odot}
  \end{align*}
  \hfill(2 points)

  For maximum frequency, we have to take highest white dwarf mass.
  \hfill(2 points)
  \begin{align*}
    f_{\mathrm{WD,max}}
      &= 2.600 \times 10^{-6} \times
         \sqrt{\frac{1.989 \times 10^{30}}{(6000 \times 10^{3})^{3}}}
         \times 1.44 \\
      &= 2.600 \times 10^{-6} \times 95.96 \times 10^{3} \times 1.44 \\
    f_{\mathrm{WD,max}} &= 0.359\,\mathrm{Hz}
  \end{align*}
  \hfill(2 points)

\item[(T11.5)]
  \[ f_{grav} = \frac{1}{\pi}\sqrt{\frac{G M_{\mathrm{NS}}}{R_{\mathrm{NS}}^3}} \]
  The lowest possible frequency is when $M_{\mathrm{NS}}$ is lowest and
  $R_{\mathrm{NS}}$ is the highest. \hfill(2 points)

  For $M_{\mathrm{NS}} = M_\odot$ and $R = 15$\,km, we get
  \[ f_{\mathrm{NS,min}} = 1.996\,\mathrm{kHz} \]
  \hfill(2 points)

  Similarly, the largest possible frequency is when $M_{\mathrm{NS}}$ is the
  largest and $R_{\mathrm{NS}}$ is the smallest. \hfill(2 points)

  For $M_{\mathrm{NS}} = 3M_\odot$ and $R = 10$\,km, we get
  \[ f_{\mathrm{NS,max}} = 6.352\,\mathrm{kHz} \]
  \hfill(2 points)

\item[(T11.6)] For black holes, we have to consider $R_{ISCO}$.
  \hfill(2 points)

  Hence the equation will be,
  \[ f_{BH} = \frac{1}{\pi} \times
     \sqrt{\frac{G M_\odot}{27 R_{sch-\odot}^3}} \times
     \frac{M_\odot}{M_{BH}} \]
  \hfill(2 points)
  \[ f_{BH} = 4.396\,\mathrm{kHz} \times \frac{M_\odot}{M_{BH}} \]
  \hfill(3 points)

\item[(T11.7)] We found the frequency of the LIGO-detected wave to be
  166.67\,Hz just before merger.
  % sic: the official solution says 166.67 Hz here, although part (T11.1)
  % found 142.86 Hz (and 142.86 is used in the mass estimate below).
  As per our analysis above, only black holes can lead to emission in this
  frequency range: Black Hole. \hfill(2 points)

  By using corresponding expression,
  \[ M_{\mathrm{obj}} = \frac{4396}{142.86} M_\odot \approx 31 M_\odot \]
  \hfill(3 points)
\end{description}
